\section{The current arithmetic route and its next calculation}
\label{sec:research-state}
The source chain now supplies an actual \(\mathbb C[t]\)-module
cohomology \(Q\), a finite invariant packet inside that cohomology,
a specified nonzero extension class, an original theta boundary
for its representative defect, and an exact arithmetic formula
for the relative weight form. The local continuation makes that
last formula an explicit confluent interpolation problem.
These constructions are logically prior to any purity conclusion.
The subsequent continuation now supplies certified entries of the
original arithmetic moment matrix, the exact finite-compression
and infinite-tail resolvent comparison, and a convergence rate for
the initial resolvent including its complete off-line jets.

\subsection{A single chain of exact comparison maps}
For each actual finite packet, retaining every full order, the
maps proved in this reader are
\[
 E=\mathbb C[t]/h
 \xrightarrow{\,R_{\mathrm{ref}}\,}\mathscr B
 \xrightarrow{\,q\,}Q,\qquad
 J_ZR_{\mathrm{ref}}=1,\qquad
 DR_{\mathrm{ref}}-R_{\mathrm{ref}}A
       =\Theta\phi_*\ell .
\]
Every admitted polynomial correction keeps these same quotient
and jet maps. At its distinguished orthogonal level let \(B_{N,u}(t)=(p_{N,u}(t)-R_Z(u)(t))/h(t)\), a polynomial of degree at most \(m\). The complete calculation is
\[
\begin{gathered}
 d\nu_Z(t)=\left|\frac{g(1/2+it)}{h(1/2+it)}\right|^2
       \frac{dt}{2\pi},\quad N=d+m,\quad
 C_NP=j_hP,\quad
 K_N=C_N(H_N^\nu)^{-1}C_N^*,\\
 p_{N,u}=(H_N^\nu)^{-1}C_N^*K_N^{-1}U_\varepsilon u,\qquad
 R_mu=R_{\mathrm{ref}}u+
       \Theta(B_{N,u}(D)\phi_*),\\
 G_m=U_\varepsilon^*K_N^{-1}U_\varepsilon,\qquad
 W_m=U_\varepsilon^*K_N^{-1}
     (AK_N+K_NA^*-K_N)K_N^{-1}U_\varepsilon .
\end{gathered}
\tag{R1}
\]
The middle line evaluates the specified polynomial at the original \(D\); every source vector lies in \(V\).

Write \(v_n=j_hq_n\) and \(\kappa_n=\|q_n\|_{\nu_Z}^2\).
The original arithmetic measure now determines the exact quantity
\[
 \epsilon_m=
 \frac{|\operatorname{Re}(v_N^*K_N^{-1}v_{N+1})|
 +\sqrt{(v_{N+1}^*K_N^{-1}v_{N+1})(v_N^*K_N^{-1}v_N)
       -(\operatorname{Im}(v_N^*K_N^{-1}v_{N+1}))^2}}
      {\kappa_N}.
 \tag{R2}
\]
The proofs in Section~\ref{sec:tau-boundary-kernel} establish
(R1)--(R2) before passing to cohomology.
For a reflection-stable packet the first absolute-value term is
zero by the exact trace identity. The same sequence obeys
\[
 \epsilon_m\ \geq\
       \sum_{\rho\in Z}m_\rho
                   |\operatorname{Re}\rho-\tfrac12|.
 \tag{R3}
\]
Thus the rank-two boundary detects the combined displacement of
the entire packet, not only its largest individual displacement.
The estimate concerns the distinguished \(R_m\) sequence.
The larger family of all polynomial source corrections has the
separately proved metric image and sharp infimum in
Section~\ref{sec:tau-chain-new}; those results do not identify
its arbitrary member with \(R_m\).

The full two-row source observation in
Section~\ref{sec:tau-boundary-kernel} is injective after the
original \(A\)-iterates are retained. It therefore offers a
finite arithmetic calculation preserving every nilpotent jet:
compute the source moments, the adjacent monic polynomials and
their complete residue vectors, rather than discard a component
whose scalar trace vanishes.

\subsection{Topology, duality and the two sides of a weight estimate}
The actual source moments now have the absolutely convergent
incomplete-gamma expression and explicit double-tail bound of
Section~\ref{sec:arithmetic-moments}. Its certified six-by-six
Gram matrix concerns the original \(|2\xi|^2dt/(2\pi)\) measure.
The exact identity (AM.19) supplies the original monic recurrence
norms through degree five, with interval enclosures.

For the \(|g/h|^2\) kernel, Section~\ref{sec:kernel-resolvent}
proves a universal polynomial quadrature formula, including
centres that coincide with a quadrature node. On the precise
nonintersection domain its alternate expression is
\[
 K_{n-1}=\kappa_{n-1}^{-1}q_n(A)S_nq_n(A)^*,\qquad
 AS_n+S_nA^*-S_n=cr_n^*+r_nc^*,\quad r_n=F_n(A)c.
\]
The exact Schur comparison (KR.15)--(KR.20) then joins the initial
and terminal observations to the actual infinite-dimensional tail.
Section~\ref{sec:gaussian-resolvent-convergence} proves the explicit
initial-resolvent rate (GC.8) and its complete finite-algebra
bound (GC.23). The projection to that off-line domain retains
the critical summand as its stated kernel. All factors in the
Schur identity remain in converting the initial estimate to the
terminal form that determines \(\epsilon_m\).

In particular, Section~\ref{sec:metric-departure} identifies
exactly the two contributions in a reflection-stable packet:
\[
 \epsilon_m^2=2\sum_{\rho\in Z}m_\rho
       (\operatorname{Re}\rho-1/2)^2+\mathfrak d_{G_m}(A).
\]
Its full-jet inequalities also quantify the required metric
ratios along each retained nilpotent chain. These identities
do not infer a relative estimate from absolute matrix shrinkage.

The theta-polynomial span is dense in every finite differential
graph norm. The completed jet graph, however, is
\(\mathscr H^{(r)}\oplus E\), with quotient \(E\), not a
quotient of \(\mathscr H^{(r)}\) in which the arithmetic
component has vanished. This supplies the exact comparison
between the shrinking Hilbert representative and its unchanged
cohomological class.

The exponential observation norm supplies a second exact
comparison. A Mellin jet at \(\rho\) is bounded precisely in
the strip
\[
 |\operatorname{Re}\rho-\tfrac12|<a,
\]
with the factorial norm in Theorem~\ref{thm:weighted-jet}.
The original differential then has the energy term
\(2a\,\operatorname{sgn}(\log x)\) in
Theorem~\ref{thm:weighted-energy}.
Consequently stronger topology and weight control are related
by that explicit term; choosing a positive observation norm
does not remove it.

For any packet \(Z\), reflection gives the conjugate-linear ring
isomorphism
\[
 \dagger_Z:E_Z\longrightarrow E_{Z^\dagger},\qquad
 f^\dagger(s)=\overline{f(1-\overline s)},\quad
 Z^\dagger=\{1-\overline\rho:\rho\in Z\},
\]
with all multiplicities retained. In fact
\(h_Z^\dagger=(-1)^dh_{Z^\dagger}\); this proves well-definedness,
and applying the same operation twice is the identity. Restrict
the following internal residue pairing to reflection-stable
\(Z\), so that this map is an involution on \(E_Z\).
The arithmetic residue form is the source form
\[
 \mathcal R_Z(f,u)=\sum_{\rho\in Z}
             \operatorname{Res}_\rho\frac{f^\dagger u}{g}\,ds .
\]
Its Jacobian contraction is
\[
 \mathcal R_Z(f,g'u)
 =\sum_{\rho\in Z}m_\rho\overline{f(1-\overline\rho)}u(\rho).
 \tag{R4}
\]
Indeed locally \(g=z^{m_\rho}w(z)\) gives
\(g'/g=m_\rho/z+w'/w\), proving (R4) with multiplicities.
The trace form pairs reflected centres and can retain
nontrivial null directions from nilpotents.
In the retained basis \(1,t,\ldots,t^{d-1}\), define
\[
 S_{ij}=\sum_{\rho\in Z}\operatorname{Res}_\rho
       \frac{(s^i)^\dagger s^j}{g(s)}\,ds,
 \quad 0\leq i,j<d,\qquad J_g=M_{j_hg'}.
\]
Changing representatives by \(h\) makes the local quotient by
\(g\) holomorphic, proving that this matrix represents the
intrinsic form. By (R4) the trace pairing has matrix \(SJ_g\).
Its exact comparison to the positive source metric is therefore
\(G_m^{-1}SJ_g\), since
\(f^*G_m(G_m^{-1}SJ_g)u=f^*SJ_gu\).
Its positivity is an arithmetic question about this operator,
not a consequence of the positive Gram construction.

The finite-field two-sided weight theorem remains a source for
methods, with its compact-support comparison, duality and
tensor hypotheses. This reader has not constructed a realization
meeting those geometric hypotheses for the entire arithmetic
quotient. It has instead calculated the current characteristic-zero
objects and their original maps, including the local-cohomology
spectral extraction and meromorphic complement documented in
the complete finite-comparison appendix.

\subsection{Provenance and mathematical status}
The coherent interpolation now has the full coefficient identity
\(J_{h,N}=U_{\upsilon_h}C_N\), with
\(\upsilon_h=j_h(g/h)\) and \(\varepsilon_h=\upsilon_h^{-1}\).
Equations (CT.3)--(CT.7) prove equality with (R1), including the
original representative, and compatibility with the complete CRT
injections between finite packets.  The actual reflection preserves
the unique minimum and proves the same-metric signed pairing of
(CT.10) and (CT.16a).

On the original tensor complex, the complete calculation is
\[
 \begin{gathered}
 E^{(k)}=E_h^{\otimes k},\qquad A^{(k)}=\sum_{j=1}^kA_{h,j},\qquad
 \mathcal I_{h,k,M}=\mathcal P_M^{(k)}
                     \cap(h(s_1),\ldots,h(s_k)),\\
 R_M=(I-P_{\mathcal B_M})s_h^{\otimes k},\quad
 G_M=(J_{k,M}M_{k,M}^{-1}J_{k,M}^*)^{-1},\\
 W_M=-(Y_M^*C_M+C_M^*Y_M),\quad
 G_M-G_{M+1}=Y_M^*Y_M,\quad
 Y_M,C_M:E^{(k)}\longrightarrow\mathcal E_M,\\
 \dim\mathcal E_M=\binom{M+k}{k-1}.
 \end{gathered}
 \tag{R5}
\]
Section~\ref{sec:coherent-tensor-integration} proves the original
ordered division, both cochain signs, every layer map, and the
extra orthogonal projection required when tensor packets are
enlarged.  Its formula (CT.18) calculates the resulting change in
both \(G\) and \(W\).

For a reflection-stable packet, let \(p_M\) be either signed
inertia of \(W_M\), and let \(\mathfrak d_M^{(k)}\) be its
complete tensor departure from (CT.20).  The same source layers give
\[
 \begin{gathered}
 \mathcal D_k\le p_M\epsilon_M,\qquad
 p_M\epsilon_M^2\ge
  2kd^{k-1}\sum_{\rho\in Z}m_\rho
        (\operatorname{Re}\rho-1/2)^2+\mathfrak d_M^{(k)},\\
 p_M\le\min\{\operatorname{rank}Y_M,
              \operatorname{rank}C_M,\lfloor d^k/2\rfloor\},\qquad
 \mathcal D_k^2\le
 \operatorname{Tr}(G_M^{-1}(G_M-G_{M+1}))
 \operatorname{Tr}(G_M^{-1}C_M^*C_M).
 \end{gathered}
 \tag{R6}
\]
The zero-inertia case is proved directly in (CT.25).  The full
ordered sum \(\mathcal D_k\) and all multiplicities are retained
in (CT.20).  The joint density and exact graph closure in
Section~\ref{sec:joint-density-dissipation} show
\(G_M\downarrow0\) for this actual total-degree sequence and
identify every retained jet in the vertical fibre of (JD.11).
The same section proves the coercive inverse-interpolation limit
with its arithmetic unit, the supported-zero limit morphism, and
the accumulated determinant cost (JD.14)--(JD.15).  These formulas
connect the relative layer energy to the exact rate of determinant
loss through the same original matrices.

The September 11 support, Rees, comparison and audit packages
are retained in the companion source directory. The current
analytic baseline is the September 12 sequence: finite operator
comparison (PR~8), chain descent (PR~9), global retraction
(PR~10), homotopy control (PR~11), formal tau recovery (PR~12)
orthogonal boundary control (PR~13), and coherent interpolation
with joint tensor boundaries (PR~14).
Part~II includes the complete selected source proof notes with
their exact revisions. Historical statements that a map remains
unconstructed belong to their stated source dates; a later
explicit construction in this reader or the source chain
supersedes that particular statement.

The new local proofs are the full support-index equivalence,
the corrected carrier and sector comparisons, the arithmetic
confluent-kernel minimum and complete two-row source observation,
the classification of source-metric images and aggregate control,
and the graph/weighted-jet comparison.
The next tranche adds the certified original theta-moment series
and its full truncation bound, the confluent quadrature and
finite/terminal/tail resolvent identities, the explicit initial
convergence rate, and the metric-departure and nilpotent-chain
identities. The recovered June 24 historical packet calculation
is reproduced with a full-jet comparison in
Section~\ref{sec:historical-packet-ghost}. Its exact
\(A\)-invariant saturation is every off-axis local factor;
the quotient, involutions, source-metric transport and split-zero
endomorphism descent are proved there. The original transcript
excerpt and its response locators are retained in the source
package. This is an explicitly recorded historical recovery,
with its new arithmetic saturation proof distinguished from
the original permutation calculation.
The subsequent historical reconstruction in
Section~\ref{historical-hurwitz-reconstruction-and-the-complete-laurent-jet-bridge}
proves the inverse map from a complete local Hurwitz zero-cluster
germ to every moment of its original compact probability measure,
followed by the explicit Bernstein recovery of that measure.
The monic cluster polynomial retains arbitrary zero multiplicity.
The same section proves the complete Laurent period maps and
their collision kernels.  For two distinct primes, the joint
Laurent map is surjective onto the full finite spectral algebra;
Theorem~\ref{thm:two-prime-laurent-presentation} gives every
primary kernel and the inverse original spectral coordinate,
including its truncated logarithm on each nilpotent factor.
The factorwise Split-Zero lifts and their exact comparison with
the synchronized global lift are proved in (S4)--(S5).
Their elementary algebraic and analytic steps are written here.
No claim of global mathematical priority follows from finding
or proving an identity in this research programme.

The finite validation records distinguish exact symbolic
identities, rational calibration models, numerical quadratures
and intentionally rejected negative controls. They do not count
test methods as theorems or imply analytic certification from
an \(L^2\) numerical experiment. The supplied formalization records
are also kept separate: successful exact-head CI metadata
does not turn the unformalized analytic theta arguments into
Lean proofs. No Lean run was performed for this continuation.

The current research calculation is (R2) on the actual
\(|g/h|^2\)-moments, with its determinant update and two-row
recurrence, together with the topology constants of
Theorem~\ref{thm:weighted-jet}. No uniform arithmetic estimate
forcing (R2) to vanish has been established in this reader.
RH, GRH and a global purity theorem are not conclusions of it.
The exact maps and source constraints above specify the live
route on which tandem work can now continue.
The next estimate must propagate (GC.8) through (KR.16), retaining
\(q_n^2\), \(\kappa_n\), \(a_n\), the terminal ratio \(F_n\)
and the actual tail denominator. The arithmetic moment formulas
give an executable source for those finite data. The exact
departure identity and the full-jet metric bounds specify what
a resulting relative estimate would control in the original
packet; no component is removed by changing observation.

The next step of that propagation is now calculated in
Section~\ref{sec:terminal-continuation}.  For
\(u_n=q_n(A)c\), \(w_n=q_{n-1}(A)c\), and
\(M_n=\kappa_{n-1}^{-1}K_{n-1}^{-1}\), let
\(\beta_n=u_n^*M_nu_n\), and let \(P_n\) be the exact
\(M_n\)-orthogonal projection with kernel \(\mathbb Cu_n\)
when \(\beta_n>0\).  Put
\(p_n=P_nw_n\), \(z_n=P_nAu_n\).  The universal step is
\[
 \begin{gathered}
 M_{n+1}=a_n^{-1}\left(M_n-
      \frac{M_nu_nu_n^*M_n}{a_n+\beta_n}\right),\qquad
 \frac{\det G_{m+1}}{\det G_m}
             =\frac{a_n}{a_n+\beta_n},\\
 \Delta_{n+1}=
 \frac{\beta_n\|z_n+a_np_n\|_n^2}{a_n(a_n+\beta_n)},\qquad
 \epsilon^{(n)}=|g_Z|+\sqrt{g_Z^2+\Delta_n},\quad
 g_Z=\operatorname{Re}\operatorname{tr}A-d/2.
 \end{gathered}
 \tag{R7}
\]
Here \(m=n-d-1\).  The zero-vector exceptional step is given
completely in (KT.14), so critical centres and every polynomial
coincidence remain in (R7)'s source calculation.  The alternate
terminal Gram follows by the exact isometry (KT.15) and recurrence
(KT.16).  The full scalar-to-matrix extraction uses the doubled
off-line algebra and the explicit sequence (KT.23); at a reflected
pair its matrix coefficient is exactly
\((-1)^lF_n^{[k+l+1]}(\rho)\).  Finally (KT.24)--(KT.25)
recover these coefficients from the original initial error and
tail, with both reciprocal series and all polynomial/norm factors
displayed.  The remaining energy estimate is therefore attached to
the three explicit terms of (KT.11), including their actual cross
pairing, on the same arithmetic moment sequence.

For a packet closed under both conjugation and reflection, the
original real-coefficient calculation in (KT.28) closes this step
in terms of consecutive determinant ratios.  With the fixed
section metric \(G_{-1}\) retained, put
\(r_m=\det G_m/\det G_{m-1}\) for \(m\ge0\).
Then the actual arithmetic allowance obeys exactly
\[
 \boxed{\epsilon_m^2
   =a_{d+m+1}\frac{(1-r_m)(1-r_{m+1})}{r_{m+1}}.}
 \tag{R8}
\]
The coefficient reality and full reflection trace prove
\(\gamma_n=0\) in the original coordinates; (KT.7)--(KT.8)
then prove (R8), including zero-step cases.  This places the
relative quartet control directly on the same monic norm ratios
and consecutive Gram determinants as the certified moment problem.

This also gives a direct arithmetic bound for the original tesserine
violation channel.  Retain the original quartet element
\[
 Q(\rho)=\rho\,1+(1-\rho)\sigma
       +\overline\rho\,\kappa+(1-\overline\rho)\sigma\kappa,
 \qquad \rho=\beta+i\gamma,
\]
and its character \(c_{-+}(\rho)=2(2\beta-1)\), as calculated
in the carrier section.  If \(\gamma\ne0\), \(\beta\ne1/2\),
and \(Z\) is this entire four-point zero orbit with common full
order \(\mu\), then \(d=4\mu\) and for every \(m\ge0\),
\[
 \boxed{
 \mu^2|c_{-+}(\rho)|^2
   \le \epsilon_m^2
   =a_{4\mu+m+1}\frac{(1-r_m)(1-r_{m+1})}{r_{m+1}}.}
 \tag{R9}
\]
To prove every coefficient, conjugation and the functional equation
preserve the full order, and the four real displacements are
\(\beta-1/2,\beta-1/2,-(\beta-1/2),-(\beta-1/2)\).
The proved aggregate bound (R3) is therefore
\(\epsilon_m\ge4\mu|\beta-1/2|
=\mu|2(2\beta-1)|\).  Squaring and using (R8) gives (R9).
The sector morphism on the original supported quartet is precisely
\[
 G(\chi_{-+}):G(\mathbb C[K_4])\longrightarrow G(\mathbb C),
 \quad Q(\rho)^\bullet\longmapsto
                   c_{-+}(\rho)^\bullet,\quad\tau\longmapsto\tau.
\]
Here \(\chi_{-+}=\operatorname{ev}_{-1,+1}\) is the original group character extended
linearly; it preserves multiplication because the character does,
and hence its displayed \(G\)-lift is a semiring morphism.
At \(\beta=1/2\) this morphism takes the supported value \(e\);
the same original quartet has phase channel \(c_{--}=4i\gamma\).
The orbit then has two distinct points when \(\gamma\ne0\), so
its degree is \(2\mu\) and (R8) applies with that exact degree;
the left-hand side of the sector bound is zero.  Thus the arithmetic
volume comparison and the supported critical value use the same
typed character map, with the phase coordinate retained.

